
\par\noindent\textit{定理~\ref{thm:7-2-212} 的证明。}
仍用上述归纳方案，对 $(\ell,i)$ 作双重归纳，
证明如下引理。

\begin{lemma}{7.2.230}\label{lem:7-2-230}
考虑两个有向系统
$\{I_i\colon C_i\to C_{i+1}\}$ 与
$\{I'_i\colon C'_i\to C'_{i+1}\}$。
固定一个 $i$，取 $C'_i$ 的模型
$\widehat{\mathfrak C}'_i$，并给定两个态射
$\{f_i^{s=s_0}\colon C_i\to C'_i\}$
（$s_0=0,1$），以及它们的交换误差同伦
$\mathcal F_i^{s=s_0}$。
令 $\mathcal H_i$、$\mathcal H_{i+1}$ 与
$\mathfrak H_i$ 如定义~\ref{def:7-2-204}。
则可将 $\mathcal H_i$ 与 $\mathfrak H_i$
提升为滤过
$A_{\bar n(\lambda_{i+\ell}),\bar k(\lambda_{i+\ell})}$
同态，使 (7.2.205)、(7.2.207)
均成为该层次的滤过同态等式。
\end{lemma}

\begin{proof}
假设所有 $(\ell',i')<(\ell,i)$ 的情形均已证明，
现处理 $(\ell,i)$。记
$(n',K')=(\bar n(\lambda_{i+\ell}),
\bar k(\lambda_{i+\ell}))$。
归纳假设给出某个 $(n,K)$，
使 (7.2.205)、(7.2.207) 在 $i$ 处
已作为滤过 $A_{n,K}$ 结构的等式成立。
可假定 $((n,K),(n',K'))$ 属于
(7.2.130.1) 或 (7.2.130.2) 的类型。

注意 $f_i^{s=s_0}$、$\mathcal F_i^{s=s_0}$、
$I_i$、$I'_i$ 与 $\mathcal I'_i$
都已提升为滤过 $A_{n',K'}$ 同态。
因此 (7.2.207.2)、(7.2.207.3) 的右侧
均为滤过 $A_{n',K'}$ 同态。
取 $\beta\in\mathfrak G$、$\lVert\beta\rVert=n'$，
考虑将 $\mathfrak H_i$ 提升至滤过
$A_{n',K'}$ 同态的障碍类
\[
o_{K',\beta}(\mathfrak H_i)
\in\operatorname{Hom}\bigl(
B_{K'}\overline C_i[1],
\overline{\widehat{\mathfrak C}}'_{i+1}[1]
\bigr).
\]
其定义见 (7.2.76.2)。依其定义、函子性及
(7.2.207.2)--(7.2.207.3)，得到
\begin{align}
\overline{\operatorname{Eval}}_{r=1}\circ
o_{K',\beta}(\mathfrak H_i)
+\delta_1\bigl((\mathcal H_{i+1}\circ I_i)_{K',\beta}\bigr)
&=0,
\tag{7.2.231.1}\label{eq:7-2-231-1}\\
\overline{\operatorname{Eval}}_{s=s_0}\circ
o_{K',\beta}(\mathfrak H_i)
+\delta_1\bigl((\mathcal F_i^{s=s_0})_{K',\beta}\bigr)
&=0.
\tag{7.2.231.2}\label{eq:7-2-231-2}
\end{align}
这里 (7.2.231.1) 中的算子为
\[
\delta_1\colon
\operatorname{Hom}\bigl(B_{K'}\overline C_i[1],
\overline{\widehat{\mathfrak C}}'_{i+1}[1]\bigr)
\longrightarrow
\operatorname{Hom}\bigl(B_{K'}\overline C_i[1],
\overline{\widehat{\mathfrak C}}'_{i+1}[1]\bigr),
\]
由 $\widehat{\mathfrak m}_1$ 与
$\overline{\mathfrak m}_1$ 诱导；
(7.2.231.2) 中的 $\delta_1$ 类似。
符号 $(\mathcal H_{i+1}\circ I_i)_{K',\beta}$
指该滤过 $A_{n',K'}$ 同态的
$(K',\beta)$-分量。

由于 (7.2.207) 在 $i+1$ 处作为滤过
$A_{n',K'}$ 同态等式成立，
在 $i$ 处则作为滤过 $A_{n,K}$ 同态等式成立，
故
\begin{equation}
\overline{\operatorname{Eval}}_{s=s_0}\circ
(\mathcal H_{i+1}\circ I_i)_{K',\beta}
=\overline{\operatorname{Eval}}_{r=1}\circ
(\mathcal F_i^{s=s_0})_{K',\beta}.
\tag{7.2.232}\label{eq:7-2-232}
\end{equation}
因而三元组
\[
\bigl(
(\mathcal F_i^{s=0})_{K',\beta},
(\mathcal F_i^{s=1})_{K',\beta},
(\mathcal H_{i+1}\circ I_i)_{K',\beta}
\bigr)
\]
属于
$\operatorname{Hom}(B_{K'}\overline C_i[1],
\mathfrak K[1])$。
结合子引理~\ref{sublem:7-2-228} 与 (7.2.231)，
可取
\[
(\mathfrak H_i)'_{K',\beta}
\in\operatorname{Hom}\bigl(
B_{K'}\overline C_i[1],
\overline{\widehat{\mathfrak C}}'_{i+1}[1]
\bigr)
\]
满足
\begin{equation}
o_{K',\beta}(\mathfrak H_i)
+\delta_1\bigl((\mathfrak H_i)'_{K',\beta}\bigr)=0,
\tag{7.2.233}\label{eq:7-2-233}
\end{equation}
以及
\begin{align}
\overline{\operatorname{Eval}}_{r=1}\circ
(\mathfrak H_i)'_{K',\beta}
&=(\mathcal H_{i+1}\circ I_i)_{K',\beta},
\tag{7.2.234.1}\label{eq:7-2-234-1}\\
\overline{\operatorname{Eval}}_{s=s_0}\circ
(\mathfrak H_i)'_{K',\beta}
&=(\mathcal F_i^{s=s_0})_{K',\beta}.
\tag{7.2.234.2}\label{eq:7-2-234-2}
\end{align}
(7.2.233) 允许用 $(\mathfrak H_i)'_{K',\beta}$
把 $\mathfrak H_i$ 提升为滤过 $A_{n',K'}$ 同态；
(7.2.234) 又保证该提升满足
(7.2.207.2)--(7.2.207.3)。

接下来提升 $\mathcal H_i$，并修正
$(\mathfrak H_i)'_{K',\beta}$，
使 (7.2.207.1) 也成立。
考虑将 $\mathcal H_i$ 提升至滤过
$A_{n',K'}$ 同态的障碍类
\[
o_{K',\beta}(\mathcal H_i)
\in\operatorname{Hom}\bigl(
B_{K'}\overline C_i[1],
\overline{\widehat{\mathfrak C}}'_{i,r=1}[1]
\bigr).
\]
由 (7.2.207) 在滤过 $A_{n,K}$ 层次成立，得
\begin{equation}
\overline{\operatorname{Eval}}_{r=0}\circ
o_{K',\beta}(\mathfrak H_i)
=o_{K',\beta}(\mathcal I'_i\circ\mathcal H_i).
\tag{7.2.235.1}\label{eq:7-2-235-1}
\end{equation}
又由 (7.2.205) 在同一层次成立，得
\begin{equation}
\overline{\operatorname{Eval}}_{s=s_0}\circ
o_{K',\beta}(\mathcal H_i)
+\delta_1\bigl((f_i^{s=s_0})_{K',\beta}\bigr)=0.
\tag{7.2.235.2}\label{eq:7-2-235-2}
\end{equation}

利用 $\overline{\mathcal I}'_i$、$\overline I'_i$ 的链同伦等价性质，
可得如下子引理。

\begin{statement}{子引理}{7.2.236}\label{sublem:7-2-236}\itshape
存在
\[
\begin{aligned}
(\mathcal H_i)_{K',\beta}
&\in\operatorname{Hom}\bigl(
B_{K'}\overline C_i[1],
\overline{\widehat{\mathfrak C}}'_{i,r=1}[1]
\bigr),\\
(\operatorname{Cor1})_{K',\beta}
&\in\operatorname{Hom}\bigl(
B_{K'}\overline C_i[1],
\overline{\widehat{\mathfrak C}}'_{i+1,r=0}[1]
\bigr),
\end{aligned}
\]
使下列条件成立：
\begin{align}
o_{K',\beta}(\mathcal H_i)
+\delta_1\bigl((\mathcal H_i)_{K',\beta}\bigr)
&=0,
\tag{7.2.237.1}\label{eq:7-2-237-1}\\
\overline{\operatorname{Eval}}_{s=s_0}\circ
(\mathcal H_i)_{K',\beta}
&=(f_i^{s=s_0})_{K',\beta},
\tag{7.2.237.2}\label{eq:7-2-237-2}\\
\overline{\operatorname{Eval}}_{r=0}\circ
(\mathfrak H_i)'_{K',\beta}
&=(\mathcal I'_i\circ\mathcal H_i)_{K',\beta}
+\delta_1\bigl((\operatorname{Cor1})_{K',\beta}\bigr),
\tag{7.2.237.3}\label{eq:7-2-237-3}\\
\overline{\operatorname{Eval}}_{s=0}\circ
(\operatorname{Cor1})_{K',\beta}
&=\overline{\operatorname{Eval}}_{s=1}\circ
(\operatorname{Cor1})_{K',\beta}=0.
\tag{7.2.237.4}\label{eq:7-2-237-4}
\end{align}
\end{statement}

\par\noindent\textit{子引理~\ref{sublem:7-2-236} 的证明。}
置
\[
\begin{aligned}
\overline D_i
&=\ker\Bigl(
\overline{\operatorname{Eval}}_{s=0}\oplus
\overline{\operatorname{Eval}}_{s=1}\colon
\overline{\widehat{\mathfrak C}}'_{i,r=1}
\to\overline C'_{i,s=0,r=1}\oplus
\overline C'_{i,s=1,r=1}
\Bigr),\\
\overline D_{i+1}
&=\ker\Bigl(
\overline{\operatorname{Eval}}_{s=0}\oplus
\overline{\operatorname{Eval}}_{s=1}\colon
\overline{\widehat{\mathfrak C}}'_{i+1,r=0}
\to\overline C'_{i+1,s=0,r=0}\oplus
\overline C'_{i+1,s=1,r=0}
\Bigr).
\end{aligned}
\]
它们是滤过 $A_{n,K}$ 代数，
$\mathcal I'_i$ 诱导链同伦等价
\begin{equation}
(\overline{\mathcal I'_i})_1\colon
\overline D_i\longrightarrow\overline D_{i+1}.
\tag{7.2.238}\label{eq:7-2-238}
\end{equation}
由 (7.2.175.3)，可取
$(\mathcal H_i)^I_{K',\beta}$
满足 (7.2.237.2)，但未必满足
(7.2.237.1)、(7.2.237.3) 或 (7.2.237.4)。
因而
\begin{equation}
o_{K',\beta}(\mathcal H_i)
+\delta_1\bigl((\mathcal H_i)^I_{K',\beta}\bigr)
\in\operatorname{Hom}\bigl(
B_{K'}\overline C_i[1],\overline D_i[1]
\bigr).
\tag{7.2.239}\label{eq:7-2-239}
\end{equation}
由 (7.2.233)、(7.2.235.1)，得到
\[
\begin{aligned}
(\overline{\mathcal I'_i})_1\circ
\bigl(o_{K',\beta}(\mathcal H_i)
+\delta_1((\mathcal H_i)^I_{K',\beta})\bigr)
=\delta_1\Bigl(
(\mathcal I'_i\circ\mathcal H_i)_{K',\beta}
-\overline{\operatorname{Eval}}_{r=0}\circ
(\mathfrak H_i)'_{K',\beta}
\Bigr).
\end{aligned}
\]
由 (7.2.206)、(7.2.237.2)，有
\[
\overline{\operatorname{Eval}}_{s=s_0}\circ
(\mathcal I'_i\circ\mathcal H_i^I)_{K',\beta}
=(I'_i\circ f_i^{s=s_0})_{K',\beta}.
\]
另一方面，依 (7.2.234.2)、(7.2.186)，
\[
\begin{aligned}
\overline{\operatorname{Eval}}_{s=s_0}\circ
\overline{\operatorname{Eval}}_{r=0}\circ
(\mathfrak H_i)'_{K',\beta}
&=\overline{\operatorname{Eval}}_{r=0}\circ
(\mathcal F_i^{s=s_0})_{K',\beta}\\
&=(I'_i\circ f_i^{s=s_0})_{K',\beta}.
\end{aligned}
\]
由于 (7.2.238) 是链同伦等价，
可取
\[
(\operatorname{Cor2})_{K',\beta}
\in\operatorname{Hom}\bigl(
B_{K'}\overline C_i[1],\overline D_i[1]
\bigr)
\]
使
\[
o_{K',\beta}(\mathcal H_i)
+\delta_1\bigl((\mathcal H_i)^I_{K',\beta}\bigr)
=\delta_1\bigl((\operatorname{Cor2})_{K',\beta}\bigr).
\]
置
\[
(\mathcal H_i)^{II}_{K',\beta}
=(\mathcal H_i)^I_{K',\beta}
-(\operatorname{Cor2})_{K',\beta}.
\]
它满足 (7.2.237.1)、(7.2.237.2)。

下面继续修正 $(\mathcal H_i)^{II}_{K',\beta}$，
并选取 $(\operatorname{Cor1})_{K',\beta}$，
使 (7.2.237.3)--(7.2.237.4) 也成立。
由于 (7.2.237.1) 已成立，
可用 $(\mathcal H_i)^{II}_{K',\beta}$
把 $\mathcal H_i$ 提升为滤过 $A_{n',K'}$ 同态，
记作 $\mathcal H_i^{II}$。
由 (7.2.233)、(7.2.235.1) 与 (7.2.237.1)，
有
\[
\begin{aligned}
\delta_1\bigl(
\overline{\operatorname{Eval}}_{r=0}\circ
(\mathfrak H_i)'_{K',\beta}\bigr)
&=-\overline{\operatorname{Eval}}_{r=0}\circ
o_{K',\beta}(\mathfrak H_i)\\
&=-o_{K',\beta}(\mathcal I'_i\circ\mathcal H_i)
=\delta_1\bigl(
(\mathcal I'_i\circ\mathcal H_i^{II})_{K',\beta}
\bigr).
\end{aligned}
\]
因此
\begin{equation}
\delta_1\Bigl(
\overline{\operatorname{Eval}}_{r=0}\circ
(\mathfrak H_i)'_{K',\beta}
-(\mathcal I'_i\circ\mathcal H_i^{II})_{K',\beta}
\Bigr)=0.
\tag{7.2.240}\label{eq:7-2-240}
\end{equation}
又由 (7.2.234.2)、(7.2.186)，
\[
\begin{aligned}
\overline{\operatorname{Eval}}_{s=s_0}\circ
\overline{\operatorname{Eval}}_{r=0}\circ
(\mathfrak H_i)'_{K',\beta}
&=\overline{\operatorname{Eval}}_{r=0}\circ
(\mathcal F_i^{s=s_0})_{K',\beta}\\
&=(I'_i\circ f_i^{s=s_0})_{K',\beta}.
\end{aligned}
\]
另一方面，依 (7.2.206)、(7.2.237.2)，
\[
\begin{aligned}
\overline{\operatorname{Eval}}_{s=s_0}\circ
(\mathcal I'_i\circ\mathcal H_i^{II})_{K',\beta}
&=(I'_i\circ
\overline{\operatorname{Eval}}_{s=s_0}\circ
\mathcal H_i^{II})_{K',\beta}\\
&=(I'_i\circ f_i^{s=s_0})_{K',\beta}.
\end{aligned}
\]
由上述端面等式可得\translatornote{原书式 (7.2.241) 的复合左因子
写作 $\mathcal I_i'^{(r=0)}$，
但本段前后都用已定义的
$\mathcal I'_i\colon
\widehat{\mathfrak C}'_{i,r=1}\to
\widehat{\mathfrak C}'_{i+1,r=0}$，
并未另定义上标 $(r=0)$；
此处按既定映射记号排印。}：
\begin{equation}
\begin{aligned}
\overline{\operatorname{Eval}}_{r=0}\circ
(\mathfrak H_i)'_{K',\beta}
-(\mathcal I'_i\circ\mathcal H_i^{II})_{K',\beta}
\in\operatorname{Hom}\bigl(
B_{K'}\overline C_i[1],\overline D_{i+1}[1]
\bigr).
\end{aligned}
\tag{7.2.241}\label{eq:7-2-241}
\end{equation}
利用 (7.2.238) 的滤过 $A_{n,K}$ 同伦等价，
可取\translatornote{原书此处把
$(\operatorname{Cor1})_{K',\beta}$
放在未移位的
$\operatorname{Hom}(B_{K'}\overline C_i,\overline D_{i+1})$；
它在子引理 (7.2.237) 与本段的
$\delta_1$ 方程中均须使用移位后的源靶，
此处按同一方程定型。}
\[
\begin{aligned}
(\operatorname{Cor1})_{K',\beta}
&\in\operatorname{Hom}\bigl(
B_{K'}\overline C_i[1],\overline D_{i+1}[1]
\bigr),\\
(\operatorname{Cor3})_{K',\beta}
&\in\operatorname{Hom}\bigl(
B_{K'}\overline C_i[1],\overline D_i[1]
\bigr),
\end{aligned}
\]
使 $\delta_1((\operatorname{Cor3})_{K',\beta})=0$，并且
\[
\begin{aligned}
\overline{\operatorname{Eval}}_{r=0}\circ
(\mathfrak H_i)'_{K',\beta}
-(\mathcal I'_i\circ\mathcal H_i^{II})_{K',\beta}
-(\overline{\mathcal I'_i})_1\circ
(\operatorname{Cor3})_{K',\beta}
=\delta_1\bigl((\operatorname{Cor1})_{K',\beta}\bigr).
\end{aligned}
\]
现在令
\[
(\mathcal H_i)_{K',\beta}
=(\mathcal H_i)^{II}_{K',\beta}
+(\operatorname{Cor3})_{K',\beta}.
\]
直接核对可知它具有所需性质，
子引理~\ref{sublem:7-2-236} 得证。

再看 (7.2.237.4)。它说明四元组
\[
\bigl(0,0,(\operatorname{Cor1})_{K',\beta},0\bigr)
\]
属于
\[
\begin{aligned}
&\bigoplus_{s_0=0}^{1}\operatorname{Hom}\bigl(
B_{K'}\overline C_i[1],
\overline{\widehat{\mathfrak C}}'_{i+1,s=s_0}[1]
\bigr)\\
&\qquad\oplus
\bigoplus_{r_0=0}^{1}\operatorname{Hom}\bigl(
B_{K'}\overline C_i[1],
\overline{\widehat{\mathfrak C}}'_{i+1,r=r_0}[1]
\bigr),
\end{aligned}
\]
且
\[
\overline B\circ
\bigl(0,0,(\operatorname{Cor1})_{K',\beta},0\bigr)=0.
\]
这里 $\overline B$ 由图式~\ref{diag:7-2-6}
中的 $B$ 诱导。该图式的正合性
给出
\[
(\operatorname{Cor4})_{K',\beta}
\in\operatorname{Hom}\bigl(
B_{K'}\overline C_i[1],
\overline{\widehat{\mathfrak C}}'_{i+1}[1]
\bigr),
\]
满足
\[
\begin{aligned}
\overline{\operatorname{Eval}}_{r=0}\circ
(\operatorname{Cor4})_{K',\beta}
&=(\operatorname{Cor1})_{K',\beta},\\
\overline{\operatorname{Eval}}_{r=1}\circ
(\operatorname{Cor4})_{K',\beta}
&=\overline{\operatorname{Eval}}_{s=0}\circ
(\operatorname{Cor4})_{K',\beta}
=\overline{\operatorname{Eval}}_{s=1}\circ
(\operatorname{Cor4})_{K',\beta}=0.
\end{aligned}
\]
最后对交换误差同伦作如下修正：
\begin{equation}
(\mathfrak H_i)_{K',\beta}
=(\mathfrak H_i)'_{K',\beta}
-\delta_1\bigl((\operatorname{Cor4})_{K',\beta}\bigr).
\tag{7.2.242}\label{eq:7-2-242}
\end{equation}
直接核对可知，这个分量与子引理~\ref{sublem:7-2-236}
中的 $(\mathcal H_i)_{K',\beta}$
具有所需性质。引理~\ref{lem:7-2-230} 得证。
\end{proof}

定理~\ref{thm:7-2-212}
由引理~\ref{lem:7-2-230} 得出。
